\chapter[Conclusion]{Conclusion}
\label{cp:Conclusion}

\parindent0pt

\section{Project Summary}

We tested one claim: \emph{the filtering that makes a multi-agent fleet resilient also makes it ignorant}. VeriFly is a
decentralised search-and-rescue simulator in which drones map unknown 3D terrain from their own cameras, divide the work
through a geodesic Voronoi partition, gossip over a lossy radio, and cope with one or two teammates who lie about where
they are and what they have searched. Four fleets run the same program, differ only in what they trust, and are paired
down to the last random number.

Resilient consensus does what it promises: W-MSR keeps a lying observer from dragging a survivor's position estimate
(0.2--0.7\,m error, against the full size of the lie for the mean), and the 2D studies confirm the robustness theory
behind it. But it buys \emph{nothing} in the search. A fleet that filters numbers but believes claims misses one
survivor in six, the same as a fleet with no defence, and almost half with two liars, because ``nobody here'' is not a
number.

The gated fleet, which separates \emph{knowing} from \emph{believing}, found \textbf{every survivor in all 140 sweep
runs} and in all 30 E1 districts. It caught every liar that faked its position, usually within seconds, and never
accused an honest drone, as guaranteed by Proposition~\ref{prop:soundness} and confirmed in every server run. The price
is about 90\,s at the baseline, mostly paid even when nobody lies. An oracle fleet told the liar from the start is no
faster, so the fleet pays for caution about \emph{where} other drones might be, not for uncertainty about \emph{who}
is lying.

\begin{keybox}[Takeaways]
\begin{enumerate}
    \item \textbf{Trust costs survivors}. Once both fleets share a predictive safety check, collisions are rare for
          both, with one exception: a hovering liar that claims to be far away (E2), which motivates a full shell check
          in the gate.
    \item \textbf{Knowledge is cheap to compute}: an exact possible-worlds model over all sets of at most $F$ liars is a
          few bitwise operations per claim, and provenance bitmasks make ``undo the liar'' automatic.
    \item \textbf{Sound evidence comes from physics}: a liar cannot fake a time-of-flight range, an empty camera view or
          a found survivor. Accusations alone prove only ``one of us two lies'', which makes framing impossible.
    \item \textbf{Caution is a dial}: the audit fraction $q$, the gate horizon and the fleet size trade time for
          certainty, and the resulting costs are measurable.
\end{enumerate}
\end{keybox}

\section{Future Outlook}

\begin{itemize}
    \item \textbf{A complete gate.} Check the \emph{whole} reachable part of a proven liar's range shell, not only
          the direction of motion, and keep a proven liar's last located position as an obstacle. This removes the
          collisions seen in E2 against a hovering liar that claims to be far away.
    \item \textbf{Cheaper caution.} Most of the time cost is present without any liar. A gate that reasons over the
          full shell geometry instead of holding conservatively, and that schedules audits during exploration instead
          of after it, should recover much of the 90\,s.
    \item \textbf{Stronger proofs.} Rule R7 (a coverage claim that is impossible from where the drone was, checked by
          range-based localisation of the liar) would catch coverage-only liars deterministically instead of
          probabilistically.
    \item \textbf{More realistic sensing.} Robust versions of R1 for heavy-tailed UWB noise, drone bodies with a radius
          and a terrain term in the gate, and wind would bring the model closer to real hardware.
    \item \textbf{Smarter adversaries.} Coordinated liars ($F=2$ colluding), liars that equivocate over unsigned
          channels, and adaptive liars that learn the audit pattern.
    \item \textbf{Hardware.} The agent code depends only on a drone's own observation, so it could run onboard
          small multirotors with UWB radios. A small indoor testbed with three or four drones is a natural next step.
\end{itemize}
